\documentclass[12pt]{article}
\usepackage[utf8]{inputenc}
\usepackage[T1]{fontenc}
\usepackage{hyperref}
\title{Tetserret Language Resources for K-12 Curriculum Development}
\author{Compiled by Claude Sonnet 5 (claude-sonnet-5) for Andr\'as Kornai}
\date{}
\begin{document}
\maketitle

\section{Introduction}
Tetserret (ISO 639-3 tez; Glottocode tets1235) is a Western Berber language of Niger. It is spoken by the Ait-Awari and Kel Eghlal in the Akoubounou commune, and there are about 2,000 speakers. UNESCO classifies it as severely endangered, and it has no established writing system. Overview: \url{https://en.wikipedia.org/wiki/Tetserret_language}.

\section{Primary Educational Resources}
None found.

\section{Linguistic and Reference Materials}
The main source is Lux (2011), a French doctoral dissertation from Lyon 2, \emph{\'Etude descriptive et comparative d'une langue menac\'ee : le tetserret, langue berb\`ere du Niger}. A 2013 published edition of 560 pages is reported to contain a glossary and texts. Review by Sarvasy (2017): \url{https://researchportalplus.anu.edu.au/en/publications/review-of-le-tetserret-langue-berb\%C3\%A3re-du-niger/}.
Comparative word list entry: \url{https://asjp.clld.org/languages/TETSERRET_2}.
Related family overviews: \url{https://en.wikipedia.org/wiki/Western_Berber_languages}.

\section{Religious and Community Resources}
None found.

\section{Digital and Audio Resources}
None found beyond the corpus extracts reported in Lux's work.

\section{Curriculum Development Context}
Documentation is essentially one scholarly description in French. There is no standard orthography and no known school use. Any curriculum would require orthography development with the community, using Lux's glossary and texts as a starting point.
\end{document}